% ===========================================================================
%  comune/frontespizio-utente.tex
% ===========================================================================
\begin{titlepage}
\thispagestyle{empty}
\begin{tikzpicture}[remember picture, overlay]

  \fill[lyraInchiostro] (current page.north west)
    rectangle ([yshift = -9mm]current page.north east);

  \node[anchor = north west, inner sep = 0pt]
    at ([shift = {(28mm, -42mm)}]current page.north west)
    {\begin{minipage}{125mm}
       {\sffamily\fontsize{10}{12}\selectfont
        \color{lyraAccentoScuro}\bfseries
        LYRA CAPITAL\par}
       \vspace{3mm}
       {\sffamily\fontsize{40}{44}\selectfont\color{lyraInchiostro}
        Manuale\\utente\par}
       \vspace{5mm}
       {\color{lyraAccento}\rule{34mm}{1.6pt}\par}
       \vspace{6mm}
       {\sffamily\fontsize{12}{17}\selectfont\color{lyraInchiostro!62}
        Che cosa fa ogni schermata, in quale ordine\\
        si usano, e che cosa significano i numeri.\par}
     \end{minipage}};

  \node[anchor = south west, inner sep = 0pt]
    at ([shift = {(28mm, 34mm)}]current page.south west)
    {\begin{minipage}{125mm}
       \begin{tcolorbox}[enhanced, colback = lyraPannello,
                         colframe = lyraBordo, boxrule = 0.6pt, arc = 4pt,
                         left = 10pt, right = 10pt, top = 9pt, bottom = 9pt]
         {\sffamily\footnotesize\color{lyraInchiostro}%
          \textbf{Come si legge.} Ogni capitolo è una zona dell'applicazione.
          Si può leggere dall'inizio, oppure aprire direttamente il capitolo
          della schermata che si ha davanti.\par
          \vspace{4pt}
          \textbf{Per modificare Lyra}, non per usarla, il documento è un
          altro: la \emph{Guida dello sviluppatore}, nella stessa cartella.\par}
       \end{tcolorbox}
     \end{minipage}};

  \node[anchor = south east, inner sep = 0pt, align = right]
    at ([shift = {(-24mm, 18mm)}]current page.south east)
    {\sffamily\scriptsize\color{lyraInchiostro!45}
     Documento generato da sorgente \LaTeX\\
     \today};

\end{tikzpicture}
\end{titlepage}

\cleardoublepage

\thispagestyle{empty}
\vspace*{18mm}

{\sffamily\Large\bfseries\color{lyraInchiostro}Prima di cominciare\par}
\vspace{3mm}
{\color{lyraAccento}\rule{\textwidth}{1pt}\par}
\vspace{7mm}

\noindent
Lyra serve a costruire strategie di investimento, provarle sui dati del
passato e — quando reggono — farle girare su un conto vero. Le tre cose si
fanno in quest'ordine, e il manuale segue lo stesso ordine.

\vspace{4mm}
\noindent
Nel testo, \voce{Questa} è una voce di menu o una scheda, \pulsante{Questo} un
pulsante da premere, \campo{questo} un campo da compilare, e
\percorso{Settings \freccia Data} significa: apri \voce{Settings}, poi la
scheda \voce{Data}.

\vspace{6mm}
{\sffamily\large\bfseries\color{lyraInchiostro}I riquadri\par}
\vspace{3mm}

\begin{nota}{Nota}
Un chiarimento. Saltandolo, l'operazione riesce lo stesso.
\end{nota}

\begin{pratica}{In pratica}
L'operazione, passo per passo.
\end{pratica}

\begin{attenzione}{Attenzione}
Qui si perde qualcosa: dati, denaro, o tempo. Vale la pena leggerlo prima.
\end{attenzione}

\vspace{6mm}
\begin{perche}{Una parola sui numeri}
Nessun risultato mostrato dall'applicazione è una previsione. Un backtest dice
come una strategia si \emph{sarebbe} comportata sui dati del passato, con i
costi stimati; non dice come si comporterà. Il capitolo sugli \voce{Streams}
spiega quali numeri guardare per distinguere una strategia che ha funzionato
da una che ha avuto fortuna.
\end{perche}

\cleardoublepage
